\subsection{A fully admissible reference does not supply the prescribed truncation}

\begin{theorem}[AT28: Proved; a coherent-reference instance with no prescribed low-form completion]
\label{AT28}
There are the integral admissible theories $A,G$ and the actual
C21 partial model $f:P\to A$ of AT12, together with strict maps
\[
 g:P\longrightarrow G,\qquad r:G\longrightarrow A,\qquad
 F:P\longrightarrow A^{\Delta[1]},\qquad
 e_0F=rg,\quad e_1F=f.
 \tag{AT.196}
\]
The reference $G$ is admissible and is a simplexwise coefficient
resolution in every weight and every form degree.
The terms of $P,A,G$ are torsionfree, $P$ has coefficient $H^0$
and rational coefficient comparison in every simplex and weight,
and all these objects satisfy the support bound $q\leq w$.
Nevertheless, for every odd prime $p$ there is no admissible
$E$ over $\mathbb Z_{(p)}$, strict map $u:E\to A_{(p)}$, and
strict truncated isomorphism
\[
 \alpha:(P_{(p)})_{<2}\xrightarrow{\ \cong\ }E_{<2},
 \qquad u_{<2}\alpha=f_{(p),<2}.
 \tag{AT.197}
\]
In particular the actual reference and path data in
\textup{(AT.196)}, even with a reference exact in all forms,
do not provide the growing strict truncation agreements of C34
for this prescribed partial map.
\end{theorem}

\begin{proof}
Use exactly the algebras and map constructed in AT12.
Write $\pi:A\to G$ for its augmentation onto the first factor
and $r:G\to A$ for its coefficient section, so $\pi r=1$.
In simplex $n$, its actual fiber-product presentation is
\[
 A_n=\mathop{\prod\nolimits_{G_n}}_{0\leq i\leq j\leq k\leq n}
          (G_n\otimes_S B)
     =G_n\oplus\bigoplus_{i\leq j\leq k}
          (G_n\otimes_S B_+).
 \tag{AT.198}
\]
The displayed sum is finite in each simplex; products and
divided powers are those of the fiber product, as in AT12.
For each nondecreasing map $\alpha:[n]\to[1]$, define a map
$h_\alpha:A_n\to A_n$ by
\[
 h_\alpha\bigl(a,(v_{ijk})\bigr)
       =\bigl(a,(\alpha(i)v_{ijk})\bigr).
 \tag{AT.199}
\]
The integer $\alpha(i)$ is either zero or one.
At a coordinate where it is one this is the identity on
$G_n\otimes_S B$. At a coordinate where it is zero it is
the composite of that algebra's augmentation to $G_n$ and
its coefficient section. Both are actual strict PD cochain
maps; their coefficient parts agree. The fiber-product
universal property therefore proves that $h_\alpha$ is a
strict weight-preserving augmented PD cochain map.
In particular the use of a zero-or-one multiplier is justified
by these actual coordinate maps, without an assertion that
arbitrary scalar multiplication preserves an algebra product.

For a monotone map $\eta:[m]\to[n]$, the structure map in
\textup{(AT.198)} pulls the triple $(i,j,k)$ back to
$(\eta(i),\eta(j),\eta(k))$ and applies the specified map
on $G$. Consequently
\[
       \eta^*h_\alpha=h_{\alpha\eta}\eta^*.
\]
Indeed the multiplier in that triple coordinate on both
sides is the same integer $\alpha(\eta(i))$.
This proves every face and degeneracy identity for the entire
family, in every simplex degree.

Here is the actual map to the cotensor. For $a\in A_n$, and
for a simplex of $\Delta[n]\times\Delta[1]$ represented by
$\theta:[m]\to[n]$ and $\alpha:[m]\to[1]$, give its value as
$h_\alpha(\theta^*a)\in A_m$. The preceding identity makes
these values compatible with every simplex operator. Each
value is a PD cochain map in $a$, so the limit defining the
cotensor supplies a strict simplicial PD cochain map
$\mathcal H:A\to A^{\Delta[1]}$.
For the constant zero map $\alpha$ its value is $r\pi$;
for the constant one map its value is the identity. Hence
$e_0\mathcal H=r\pi$ and $e_1\mathcal H=1_A$.
Set $g=\pi f$ and $F=\mathcal H f$.
These formulas prove every identity in \textup{(AT.196)}.
They do not require that the separately displayed map
$P\to G$ in AT12 was chosen to equal $\pi f$.

All the stated properties of $A,G$ are proved in AT12.
The C21 source $P$ is built by actual relative free identity
disks. Pointwise, the free disk complements are torsionfree;
their degree-zero cohomology is the coefficient copy, and
over $\mathbb Q$ their monomial contractions give the
coefficient comparison in every form. These properties pass
through the filtered union in C21. Its assigned weights give
the displayed support condition. Thus the source properties
also hold for this precise prescribed map.

Every formula above is defined over $\mathbb Z$ and remains
the same formula after localization at any prime. The
nonexistence assertion \textup{(AT.197)} is exactly AT12's
proved obstruction for this $f$, which uses the actual
coefficient-syzygy class of AT11 and the specified simplex-two
coordinate projection. In particular it applies to the
present $g$ and $F$ without changing that prescribed map.
Any C34 agreement of depth at least two restricts to
\textup{(AT.197)}. Its depths therefore cannot tend to infinity
over the odd primes. This proves the assertion.
\end{proof}

This instance has the explicit admissible common source
$G\xrightarrow{r}A$ and $G\xrightarrow{1}G$.
The theorem concerns the fixed partial map and its strict
truncation agreements. It leaves the existential common-source
problem unchanged, and does not identify the more specialized
free-family input in AT26 with this example.
